\subsection{Arquitectura de componentes}

La Figura~\ref{fig:diagrama_arquitectura} presenta la arquitectura implementada y diferencia los componentes propios de Carer-Record de los sistemas hospitalarios externos.

\begin{figure}[H]
    \centering
    \resizebox{\textwidth}{!}{%
    \begin{tikzpicture}[
        component/.style={draw, rounded corners, align=center, minimum width=3cm, minimum height=1.2cm, fill=blue!8},
        service/.style={draw, rounded corners, align=center, minimum width=2.7cm, minimum height=1.1cm, fill=green!10},
        data/.style={draw, rounded corners, align=center, minimum width=2.7cm, minimum height=1.1cm, fill=violet!10},
        external/.style={draw, dashed, rounded corners, align=center, minimum width=3cm, minimum height=1.2cm, fill=gray!10},
        boundary/.style={draw=blue!50!black, thick, rounded corners, inner sep=8mm},
        flow/.style={-{Latex[length=2.5mm]}, thick}
    ]
        \node[component] (flutter) {Aplicación móvil\\Flutter};
        \node[component, right=2.2cm of flutter] (api) {API y lógica de negocio\\FastAPI cloud};
        \node[service, above right=1.1cm and 2.1cm of api] (stt) {Servicio ASR\\Transcripción};
        \node[service, below right=1.1cm and 2.1cm of api] (llm) {Servicio LLM\\Clasificación y síntesis};
        \node[data, right=2.2cm of stt] (auth) {Supabase Auth\\Identidad y sesión};
        \node[data, right=2.2cm of api] (db) {PostgreSQL\\Datos estructurados};
        \node[data, right=2.2cm of llm] (storage) {Supabase Storage\\Archivos privados};
        \node[service, right=2.2cm of db] (interop) {Adaptadores de interoperabilidad\\FHIR / HL7v2};
        \node[external, right=2.4cm of interop] (simulator) {Simulador de SIH/FCE\\Datos sintéticos};

        \node[boundary, fit=(flutter)(api)(stt)(llm)(auth)(db)(storage)(interop), label=above:{Frontera de Carer-Record}] (system) {};

        \draw[flow] (flutter) -- node[above]{Operaciones clínicas} (api);
        \draw[flow] ([yshift=3mm]flutter.east) to[bend left=25] node[above]{Credenciales} (auth.west);
        \draw[flow] (auth.west) to[bend left=18] node[below]{Sesión} ([yshift=-3mm]flutter.east);
        \draw[flow] (api) -- (stt);
        \draw[flow] (api) -- (llm);
        \draw[flow] (api) -- (auth);
        \draw[flow] (api) -- (db);
        \draw[flow] (api) -- (storage);
        \draw[flow] (simulator.west) -- node[above]{Datos sintéticos} (interop.east);
        \draw[flow] (interop.west) -- node[above]{Datos normalizados} (db.east);
    \end{tikzpicture}%
    }
    \caption{Arquitectura de componentes implementada en Carer-Record}
    \label{fig:diagrama_arquitectura}
\end{figure}

La aplicación Flutter concentra la interacción con el personal asistencial. FastAPI aplica autorización y reglas de negocio antes de acceder a los servicios administrados. Supabase Auth emite la sesión; PostgreSQL mantiene los datos estructurados; y Storage conserva archivos privados. Los servicios ASR y LLM se invocan desde la API para impedir que la aplicación móvil acceda directamente a credenciales de infraestructura. Los adaptadores de interoperabilidad reciben recursos FHIR o mensajes HL7v2 sintéticos, los normalizan y permiten comprobar el intercambio con un sistema hospitalario simulado dentro del alcance de Hito 3.

La separación por componentes favorece bajo acoplamiento, despliegue independiente de los procesos de IA y una frontera clara para aplicar seguridad. La implementación utiliza módulos lógicos coordinados por FastAPI, sin imponer un microservicio independiente para cada responsabilidad~\cite{arch_patterns_health}.
